\section{Derivados sobre activos no almacenables}
\label{sec:nonstorable_assets}

La valoración de derivados sobre activos físicos depende de si el subyacente puede
almacenarse y trasladarse entre periodos. En commodities almacenables, la relación
entre spot y forward puede apoyarse en argumentos de \textit{cash-and-carry}, en los
que intervienen costes de almacenamiento y \textit{convenience yield}
\cite{schwartz_1997}. Esta lógica se debilita cuando el objeto económico es
perecedero o no puede inventariarse. La electricidad constituye un precedente
útil: su limitada almacenabilidad hace que la estructura forward dependa en mayor
medida de expectativas sobre oferta y demanda, estacionalidad y primas por riesgo
\cite{lucia_schwartz_2002}.

Esta distinción es central en el presente TFM. El subyacente económico no es la
GPU como activo de capital, sino la \emph{GPU-hora} como servicio de capacidad
durante un intervalo temporal. El hardware puede adquirirse, mantenerse y
revenderse; una unidad de capacidad no utilizada en un periodo, en cambio, no puede
conservarse y venderse posteriormente como esa misma unidad de servicio. La
literatura sobre \textit{cloud computing} adopta una interpretación similar al
analizar recursos computacionales perecederos, pricing dinámico y mecanismos
\textit{on-demand} o \textit{spot}
\cite{xu2013dynamiccloud,li2016spotpricing,wu2019cloudpricing}.

La analogía con la electricidad es estructural, no literal. La capacidad de cómputo
presenta heterogeneidad por hardware, localización y calidad del servicio, y puede
sustituirse parcialmente entre proveedores, regiones o fechas
\cite{li2016spotpricing}. Por ello, el término ``no almacenable'' se utiliza aquí en
sentido económico: la unidad temporal de servicio no puede inventariarse, sin
afirmar que el hardware carezca de valor de almacenamiento ni que el mercado de
compute sea equivalente a un mercado eléctrico.

Esta característica afecta también a la valoración. Al no existir un spot de
GPU-hora directamente almacenable y negociable con las propiedades requeridas por
una replicación estándar, no se presupone una medida neutral al riesgo única
identificada mediante argumentos convencionales de \textit{carry}. La dinámica bajo
$\mathbb{Q}$ se interpreta, por tanto, como una \emph{ley de valoración modelizada}
dentro del experimento cuantitativo, no como una medida recuperada de una superficie
líquida de derivados sobre compute. La distinción entre dinámica física y de
valoración se formaliza en el Capítulo~\ref{chap:metodologia}.

Desde el punto de vista económico, la ausencia de un mercado completo no elimina
el riesgo de precio para un consumidor recurrente de capacidad. Si la exposición
depende del coste soportado a lo largo de varios periodos, resulta natural
representarla mediante un derivado ligado al precio medio, lo que conduce a las
opciones asiáticas.

\section{Valoración de opciones asiáticas}
\label{sec:asian_options_literature}

Las opciones asiáticas son derivados \textit{path-dependent}: su payoff depende de
una media de precios observados durante un periodo, no sólo del precio terminal
\cite{uned_mod7_2026}. El presente trabajo se limita a una call europea de strike
fijo, media aritmética y fixings discretos. Si
$0<t_1<\cdots<t_N=T$, la media terminal y el payoff son

\begin{equation}
    \bar P_T=\frac{1}{N}\sum_{i=1}^{N}P_{t_i},
    \qquad
    H_T=(\bar P_T-K)^+ .
    \label{eq:asian_call_payoff_literature}
\end{equation}

La dependencia de trayectoria implica que, una vez iniciado el periodo de
\textit{averaging}, el estado contractual debe conservar la información ya fijada
en la media. Por tanto, la valoración no puede reducirse únicamente al spot y al
tiempo hasta vencimiento; también debe incorporar la contribución acumulada de los
fixings realizados. La representación concreta de ese estado se define en el
Capítulo~\ref{chap:metodologia}.

La elección de la media aritmética introduce la principal dificultad numérica. Bajo
el entorno lognormal clásico, la media geométrica es analíticamente más tratable y
permite una valoración analítica útil como control de varianza
\cite{kemna_vorst_1990}. Turnbull y Wakeman proponen una aproximación para la
media aritmética \cite{turnbull_wakeman_1991}. La media aritmética, en cambio, es una
suma de variables lognormales correlacionadas y no dispone de una distribución
cerrada simple. De ahí la utilización histórica de métodos aproximados y numéricos:
Monte Carlo con reducción de varianza \cite{kemna_vorst_1990}, aproximaciones
analíticas \cite{turnbull_wakeman_1991}, condicionamiento sobre la media geométrica
\cite{curran1994asian} y técnicas de reducción dimensional y cotas
\cite{rogers1995asian}.

En este TFM no se comparan dichas aproximaciones. La media aritmética se mantiene
porque representa directamente el riesgo económico que se desea modelizar. Para
un usuario que compra capacidad de forma recurrente, el coste medio durante el
horizonte está relacionado con el promedio de los precios pagados. Una call sobre
ese promedio puede interpretarse, de forma estilizada, como protección frente a que
el coste medio supere un umbral $K$. El contrato es hipotético: no se presupone la
existencia actual de un mercado líquido de opciones asiáticas sobre GPU-hora.

La combinación de media aritmética, fixings discretos y estado contractual
acumulado convierte la valoración en un problema adecuado para simulación y para
el estudio de \textit{surrogate models}. El objetivo no es sólo aproximar el precio,
sino también su sensibilidad respecto al precio actual de compute.

\section{Modelos de commodities y reversión a la media}
\label{sec:commodity_mean_reversion}

El movimiento browniano geométrico (GBM) constituye el benchmark clásico en
valoración de opciones por su positividad, transición exacta y simplicidad
\cite{black_scholes_1973,merton_1973,uned_mod7_2026}. Sin embargo, no incorpora un
nivel de largo plazo hacia el que el precio tienda a retornar. En determinados
mercados de commodities, factores de oferta, demanda, capacidad productiva y
sustitución motivan modelos con reversión a la media. La literatura documenta esta
característica en distintos mercados y la incorpora explícitamente en modelos de
commodities y electricidad
\cite{bessembinder1995meanreversion,gibson_schwartz_1990,schwartz_1997,
lucia_schwartz_2002,cartea_figueroa_2005}.

Una representación parsimoniosa es el proceso de Ornstein--Uhlenbeck (OU)
\cite{uhlenbeck_ornstein_1930}. Para evitar precios negativos, el presente trabajo
aplica la reversión al logaritmo del precio, $X_t=\log P_t$:

\begin{equation}
    dX_t=\kappa(\theta-X_t)\,dt+\sigma\,dW_t,
    \qquad
    P_t=e^{X_t}>0,
    \qquad \kappa>0 .
    \label{eq:logou_literature}
\end{equation}

La desviación respecto a $\theta$ decae como $e^{-\kappa u}$ y la vida media de
un shock es $\log(2)/\kappa$. Además, la transición del OU es exacta y su varianza
condicional converge a $\sigma^2/(2\kappa)$. Estas propiedades hacen del Log-OU un
entorno especialmente conveniente para generar trayectorias sin introducir error
de discretización entre fechas de simulación. En régimen estacionario,
$\exp(\theta)$ corresponde a la mediana del precio, mientras que la media es
$\exp\{\theta+\sigma^2/(4\kappa)\}$.

La elección de un modelo de un factor responde a parsimonia, no a una afirmación
de verdad estructural. La literatura propone modelos más ricos con
\textit{convenience yield}, tipos de interés estocásticos, estacionalidad o saltos
\cite{gibson_schwartz_1990,schwartz_1997,lucia_schwartz_2002,
cartea_figueroa_2005}. Sin embargo, la historia disponible del benchmark H100 es
corta y no existe una curva líquida de derivados sobre compute que permita
identificar de forma robusta factores adicionales. El Log-OU se utiliza, por
tanto, como modelo principal parsimonioso, mientras que el GBM mantiene su papel
de benchmark conceptual.

\begin{table}[htbp]
\centering
\small
\caption{Comparación conceptual entre GBM y Log-OU}
\label{tab:gbm_logou_comparison}
\begin{tabular}{p{3.8cm}p{4.2cm}p{4.2cm}}
\toprule
\textbf{Propiedad} & \textbf{GBM} & \textbf{Log-OU} \\
\midrule
Positividad del precio & Sí & Sí, mediante $P_t=e^{X_t}$ \\
Reversión a la media & No & Sí \\
Nivel de largo plazo & No explícito & $\theta$ en log-precios \\
Persistencia de shocks & Sin fuerza de retorno & Decae como $e^{-\kappa u}$ \\
Varianza de largo plazo & Crece con el horizonte & Converge en log-precios \\
\bottomrule
\end{tabular}
\end{table}

Debe distinguirse, finalmente, entre la dinámica física
$(\kappa_P,\theta_P,\sigma_P)$ estimada sobre datos históricos y los parámetros de
valoración $(\kappa_Q,\theta_Q,\sigma_Q)$. En commodities líquidos, la estructura
de futuros puede aportar información sobre estas dinámicas y las primas de riesgo
\cite{gibson_schwartz_1990,schwartz_1997}. En compute no se dispone de una
estructura equivalente suficientemente líquida; por ello, $\mathbb Q$ vuelve a
interpretarse como una ley de valoración modelizada. Esta separación condiciona el
alcance de las conclusiones y evita identificar indebidamente parámetros históricos
con parámetros neutrales al riesgo.

\section{Monte Carlo y Randomized Quasi-Monte Carlo}
\label{sec:mc_rqmc_literature}

Monte Carlo es una herramienta estándar para valorar derivados cuando la esperanza
del payoff no admite una solución analítica, y resulta especialmente natural en
contratos dependientes de trayectoria
\cite{boyle1977montecarlo,glasserman_2004,uned_mod7_2026}. Para un estado
$\mathbf{x}$, el precio puede escribirse como

\begin{equation}
    V(\mathbf{x})
    =
    \mathbb E^{Q}\!\left[G(\mathbf Z;\mathbf{x})\right],
    \label{eq:mc_expectation_literature}
\end{equation}

donde $G$ es el payoff descontado y $\mathbf Z$ contiene las innovaciones que
generan la trayectoria futura. El estimador Monte Carlo convencional converge con
error de orden $M^{-1/2}$, lo que puede exigir un número elevado de simulaciones
cuando se necesitan etiquetas precisas para muchos estados
\cite{glasserman_2004}.

Quasi-Monte Carlo (QMC) sustituye números pseudoaleatorios independientes por
secuencias de baja discrepancia diseñadas para cubrir de forma más uniforme el
hipercubo de integración \cite{sobol_1967,caflisch_1998}. La aleatorización de esas
secuencias da lugar a Randomized Quasi-Monte Carlo (RQMC), que conserva la
estructura de baja discrepancia y permite utilizar replicaciones aleatorizadas para
estimar el error numérico \cite{owen_1997}. En problemas de valoración, esta
combinación puede mejorar la eficiencia de integración respecto a Monte Carlo
convencional, aunque la ganancia depende de la dimensión y regularidad efectiva del
integrando \cite{caflisch_1998,glasserman_2004}.

El interés de RQMC en este TFM es instrumental. El motor de simulación genera para
cada estado una etiqueta de precio y una Delta de referencia. La sensibilidad puede
estimarse mediante diferencias finitas o mediante un estimador \textit{pathwise}.
Las diferencias finitas son conceptualmente simples, pero requieren revalorar el
contrato bajo perturbaciones del input y su precisión depende del tamaño del
\textit{bump}. Cuando el payoff permite intercambiar diferenciación y esperanza, el
método pathwise estima directamente la derivada de cada trayectoria
\cite{broadie_glasserman_1996,glasserman_2004}. En el trabajo, las diferencias
finitas se utilizan como control numérico de las etiquetas pathwise, no como
generador principal.

Esta separación es importante para interpretar el experimento: MLP y DML no
aprenden un ``precio verdadero'' de mercado, sino la función inducida por el Log-OU
y aproximada mediante RQMC. La precisión de las etiquetas limita la información
disponible para el aprendizaje. Por ello se distingue entre conjuntos ordinarios de
entrenamiento/evaluación y un conjunto de referencia generado con mayor presupuesto
RQMC, utilizado para verificar que las conclusiones no dependan del ruido numérico
de las etiquetas.

\section{Machine Learning para pricing}
\label{sec:ml_pricing_literature}

Un motor de valoración por simulación puede ser preciso pero costoso cuando debe
evaluarse repetidamente sobre un gran número de estados. Las redes neuronales
permiten trasladar parte de ese coste a una fase \textit{offline}: se generan
etiquetas con el motor numérico y posteriormente se entrena una función
$\widehat V_{\theta}(\mathbf x)$ que actúa como \textit{surrogate}. Este enfoque no
sustituye la definición financiera del precio; aproxima la función producida por el
modelo de valoración.

La utilización de redes neuronales en derivados tiene antecedentes tempranos.
\textcite{hutchinson1994} emplean \textit{learning networks} para aproximar
funciones de pricing y estudiar cobertura fuera de muestra. Trabajos posteriores
han extendido el enfoque a valoración, volatilidad implícita y calibración
\cite{rufwang2020,liu2019optionpricing,horvath2021deeplearningvolatility}.
El atractivo del surrogate reside en combinar flexibilidad no lineal y evaluaciones
rápidas una vez concluido el entrenamiento.

En este TFM, el benchmark es un \textit{Multi-Layer Perceptron} entrenado de forma
supervisada con etiquetas de precio. Dado un conjunto
$\{(\mathbf x_i,V_i)\}_{i=1}^{N}$, el objetivo convencional penaliza el error entre
$V_i$ y $\widehat V_\theta(\mathbf x_i)$. Como la red es diferenciable, puede
obtenerse posteriormente una Delta mediante diferenciación automática:

\begin{equation}
    \widehat{\Delta}_{\mathrm{MLP}}
    =
    \frac{\partial \widehat V_\theta}{\partial P}.
    \label{eq:mlp_derived_delta}
\end{equation}

No obstante, una buena aproximación de los niveles no garantiza una aproximación
igualmente buena de las pendientes. La pérdida de precio no contiene una
restricción explícita sobre la forma local de la superficie. Este punto es
especialmente relevante en derivados, donde las Greeks son variables de gestión
del riesgo y no simples productos secundarios de la función de pricing. Por ello,
el MLP convencional proporciona un benchmark natural para evaluar el valor
incremental de utilizar información diferencial durante el entrenamiento.

\section{Differential Machine Learning}
\label{sec:dml_literature}

Differential Machine Learning (DML) amplía el aprendizaje supervisado incorporando,
junto a cada etiqueta de valor, información sobre sus derivadas respecto a los
inputs. Huge y Savine \parencite{huge_savine_2020} proponen esta formulación para
problemas financieros en los que el motor de simulación puede producir
sensibilidades conjuntamente con los valores.

En el caso del presente trabajo, la diferencia frente al MLP es deliberadamente
acotada. Ambos modelos aproximan la misma función de valoración, utilizan los mismos
estados y comparten la misma familia arquitectónica; DML añade supervisión sobre la
Delta. De forma esquemática,

\begin{equation}
    \text{MLP: }(\mathbf x,V),
    \qquad
    \text{DML: }(\mathbf x,V,\Delta).
\end{equation}

La función objetivo combina un término de precio y otro diferencial. En forma
genérica,

\begin{equation}
    \mathcal L_{\mathrm{DML}}
    =
    \mathcal L_V
    +
    \lambda_\Delta\,\mathcal L_\Delta,
    \qquad
    \mathcal L_\Delta
    \propto
    \left\|
        \frac{\partial \widehat V_\theta}{\partial P}
        -
        \Delta
    \right\|^2 ,
    \label{eq:dml_general_loss_literature}
\end{equation}

con el escalado y la ponderación definidos posteriormente en la metodología. La
supervisión diferencial restringe así no sólo el nivel de la superficie aprendida,
sino también su pendiente local. No equivale a una regularización convencional:
no se penaliza que la derivada sea grande, sino que se aleje de una etiqueta
diferencial de referencia \cite{huge_savine_2020}.

La principal motivación es la eficiencia informativa. Si el motor RQMC proporciona
precio y Delta para un mismo estado, cada ejemplo contiene información sobre el
nivel y la geometría local de la función. Huge y Savine
\parencite{huge_savine_2020} muestran que esta estructura puede mejorar la
aproximación de sensibilidades y la eficiencia muestral. En este TFM, sin embargo,
``eficiencia muestral'' se utiliza de forma estrictamente empírica: se comprobará si
DML alcanza una precisión determinada con menos estados dentro de la malla
experimental, sin atribuir una equivalencia universal entre una derivada y un número
fijo de observaciones adicionales.

DML introduce también costes y riesgos. La calidad del entrenamiento depende de la
calidad de las etiquetas diferenciales; un error sistemático en la Delta puede
convertirse en supervisión incorrecta. Además, precio y derivada deben escalarse de
forma coherente, la pérdida requiere diferenciación de orden superior respecto a los
parámetros de la red y el coste de entrenamiento aumenta frente al MLP
convencional. Finalmente, como cualquier surrogate, su validez está condicionada al
dominio cubierto por los datos y no garantiza extrapolación fiable
\cite{huge_savine_2020}.

La comparación MLP--DML se diseña precisamente para aislar esta diferencia de
información. El interés empírico no termina en comprobar si disminuye el error de
Delta: dado que la sensibilidad se utiliza después para cobertura, debe evaluarse si
la mejora estadística tiene una traducción económica.

\section{Cobertura y evaluación económica}
\label{sec:hedging_literature}

La Delta mide la respuesta local del valor del derivado al factor de riesgo,
$\Delta_t=\partial V_t/\partial P_t$. En el marco clásico de
Black--Scholes--Merton, una posición adecuada en el propio subyacente permite
neutralizar localmente esa exposición y constituye la base de la replicación
\cite{black_scholes_1973,merton_1973}. Sin embargo, la cobertura real introduce
fricciones que separan la precisión de la Greek del resultado económico de la
estrategia.

La primera es el rebalanceo discreto. Incluso utilizando la Delta correcta del
modelo, la posición permanece fija entre dos fechas mientras la sensibilidad de la
opción cambia. La literatura clásica muestra que esta discretización genera error
de cobertura \cite{boyle_emanuel_1980}; con costes de transacción, aumentar la
frecuencia de ajuste introduce además un compromiso entre error de replicación y
coste de negociación \cite{leland1985}. Por ello, una Delta exacta no implica una
replicación dinámica exacta en tiempo discreto.

La segunda dificultad aparece cuando el factor de riesgo no es directamente
negociable. La literatura de hedging con futuros y \textit{cross-hedging} estudia
precisamente posiciones en instrumentos relacionados y el riesgo residual derivado
de una relación imperfecta con la exposición original
\cite{ederington1979,anderson_danthine_1981}. Esta situación es relevante para
compute: la Delta respecto al precio de GPU-hora sigue siendo una medida de
sensibilidad, pero no implica que puedan negociarse directamente
$\Delta_t$ unidades de un spot almacenable.

El diseño empírico separa por ello dos preguntas. La cobertura local comprueba si
una Delta estimada reproduce mejor la respuesta del derivado ante una pequeña
perturbación del precio. La cobertura dinámica transforma esa sensibilidad en una
posición sobre un forward de compute modelizado y permite que contrato, factor e
instrumento evolucionen entre rebalanceos. Si
$J_t^G=\partial G_t/\partial P_t$ es la sensibilidad del valor actual del
instrumento de cobertura, la neutralidad local requiere

\begin{equation}
    h_t=\frac{\Delta_t}{J_t^G},
    \qquad
    \widehat h_t-h_t^*
    =
    \frac{\widehat\Delta_t-\Delta_t^*}{J_t^G}.
    \label{eq:forward_hedge_ratio_lit}
\end{equation}

Esta expresión muestra por qué el error de Delta puede amplificarse cuando
$|J_t^G|$ es pequeño. Bajo una dinámica con fuerte reversión a la media, un cambio
del spot actual puede tener un efecto atenuado sobre un precio forward de una fecha
posterior, de modo que la conversión desde Delta a unidades de hedge puede ser
sensible. El mecanismo no implica que toda observación extrema se explique por el
Jacobiano, pero justifica analizar por separado el error de la Greek y el error
terminal de la estrategia.

La evaluación económica debe considerar, por tanto, la distribución completa del
error de cobertura. RMSE, MAE y mediana absoluta capturan aspectos distintos y
deben complementarse con percentiles de cola y comparaciones trayectoria a
trayectoria. Esta distinción resulta especialmente importante cuando unas pocas
trayectorias concentran una fracción elevada del error cuadrático.

\section{Posicionamiento y aportación del TFM}
\label{sec:positioning_contribution}

La literatura revisada sitúa el trabajo en la intersección de tres líneas. La
primera estudia la capacidad informática como recurso perecedero y sus mecanismos de
pricing \cite{xu2013dynamiccloud,li2016spotpricing,wu2019cloudpricing}; la segunda
desarrolla métodos de valoración y cobertura mediante redes neuronales, culminando
en la incorporación explícita de información diferencial
\cite{hutchinson1994,huge_savine_2020}; la tercera analiza la relación entre Greeks,
rebalanceo y efectividad de cobertura
\cite{boyle_emanuel_1980,ederington1979,anderson_danthine_1981}.

El TFM conecta estas líneas mediante una comparación emparejada de MLP y DML
sobre una asiática de compute: evalúa si aprender mejor la Delta se traduce en
menor riesgo local y dinámico. La aportación reside en contrastar esa transmisión
dentro de un mismo modelo de valoración y distinguir sus límites contractuales.
